\chapter*{What this ledger is}
\addcontentsline{toc}{chapter}{What this ledger is}

This document does two things, and they are different in kind. Conflating them
is the failure mode it is built to avoid.

\paragraph{Half one --- a verdict.}
That \emph{one medium} supports the \textbf{linear} part of every force: all the
far-field effects of gravity, light, gravitomagnetism, charge and magnetism.
This is the part that can be fully derived, and it is the body of the document.
Each sector states what it \emph{needs} in order to survive; the brane and the
bulk are defined \textbf{last}, at the knit, by asking whether one medium
supplies every requirement at once.

This half is able to fail. A genuine no-go between two sectors' requirements
\emph{is} the falsification, and it would be a first-class result rather than a
failed run.

\paragraph{Half two --- an inventory.}
What is left, that can only be settled by \textbf{simulation}, stated precisely
enough to start that work. It is broad: the throat interior, the geon, the drain
law, the nonlinear brane-shear action, and the packet.

\textbf{Nothing requiring a simulation is done here.} Half two specifies the
work; it does not attempt it.

\paragraph{The seam between them is a seam, not a ceiling.}
Half one is linear response about the equilibrium state; half two is everything
nonlinear. That boundary marks where one half ends and the other begins --- it
is not a limit on the programme, and a missing nonlinear object is a half-two
\emph{row}, never a blocker in front of half one.

\paragraph{What this is not.}
A derivation of the universe from first principles. The claim is that
circumstances \emph{can exist} to support the forces we observe. Departing from
the behaviour of any real superfluid is expected, not a defect: real superfluids
are made of atoms with electromagnetic interactions, and their equations carry
that microphysics inside them. This model postulates a medium with whatever
properties the forces require --- and every postulate is a knob, counted against
the model.
